% ID: SEL-D0056-OB-E03
% ESTADO: cotejada. Segunda pasada visual y borrador compilado revisados.
% Fuente: Selectividad/sel_2011_matematicas_aplicadas/123-EXAMEN 1.pdf
% SHA-256: 69f8b84ed22f2b49074eafb95dbdce8a98114a4a2f02989b92f3f7add9355155
% Página: 2; opción B; ejercicio original 3.

Una compañía aseguradora realiza operaciones de seguros médicos y de seguros de vida. El 20\% de las operaciones corresponde a seguros médicos y el resto a seguros de vida. El porcentaje de operaciones en las que no se producen retrasos en los pagos es del 10\% en los seguros médicos y del 15\% en seguros de vida.
\begin{enumerate}[label=\alph*)]
\item \textbf{(1.5 puntos)} Halle el porcentaje de operaciones en las que no se producen retrasos en los pagos.
\item \textbf{(1 punto)} De las operaciones que han sufrido retrasos en los pagos, \textquestiondown{}qué porcentaje corresponde a los seguros de vida?
\end{enumerate}
